% Sampul dan lembar persetujuan pembimbing (ukuran sama dengan proposal)
% ============================ SAMPUL ============================
\thispagestyle{empty}
\begin{ruang}{0.4}\end{ruang}
\begin{sampulbesar}
\textbf{SEMINAR HASIL}
\end{sampulbesar}
\begin{ruang}{0.5}\end{ruang}
\begin{sampulbesar}
\textbf{ANALISIS KINERJA ALGORITMA PENJADWALAN PAKET PADA JARINGAN LTE MENGGUNAKAN SIMULATOR NS-3}
\end{sampulbesar}
% ruang kosong setinggi baris "Diajukan untuk memenuhi ..." pada sampul Skripsi, agar posisi semua sampul sama
\begin{ruang}{1.87}\end{ruang}
\begin{ruang}{0.6}\end{ruang}
\begin{tengah}
\logo{4.5cm}
\end{tengah}
\begin{ruang}{0.6}\end{ruang}
\begin{sampulsedang}
\textbf{Disusun oleh:}

\textbf{LA ODE CONTOH PRATAMA}

\textbf{NPM 2201010001}
\end{sampulsedang}
\begin{ruang}{0.8}\end{ruang}
\begin{sampulsedang}
\textbf{PROGRAM STUDI REKAYASA SISTEM KOMPUTER}

\textbf{FAKULTAS TEKNIK}

\textbf{UNIVERSITAS MUHAMMADIYAH BUTON}

\textbf{BAUBAU}

\textbf{2027}
\end{sampulsedang}
\begin{halamanbaru}\end{halamanbaru}

% ============================ LEMBAR PERSETUJUAN ============================
\thispagestyle{empty}
\begin{tengah}
\textbf{\underline{SEMINAR HASIL}}
\end{tengah}

\begin{tanpagaris}
\begin{longtable}{@{}p{0.16\linewidth}p{0.04\linewidth}p{0.74\linewidth}@{}}
Nama & : & La Ode Contoh Pratama \\[6pt]
NPM & : & 2201010001 \\[6pt]
\textit{Contact} & : & Jl. Contoh No. 1, Baubau \\[6pt]
 & & 0812-0000-0000; contoh@example.ac.id \\
 & & {\footnotesize (alamat, telepon, email untuk kemudahan menghubungi; dan jika ada perubahan data harap segera menghubungi Program Studi)} \\
\end{longtable}
\end{tanpagaris}

\begin{ruang}{1.5}\end{ruang}
\begin{tengah}
\textbf{PERSETUJUAN PEMBIMBING}

Laporan Seminar Hasil ini telah disetujui oleh:
\end{tengah}

\begin{tanpagaris}\begin{tabelpusat}
\begin{longtable}{@{}>{\centering\arraybackslash}p{0.48\linewidth}>{\centering\arraybackslash}p{0.48\linewidth}@{}}
\textbf{Pembimbing I}\newline ~\newline ~\newline ~\newline ~\newline ~\newline ~\newline ~\newline \satubaris{\PembimbingI}\newline NIDN. \PembimbingInidn &
\textbf{Pembimbing II}\newline ~\newline ~\newline ~\newline ~\newline ~\newline ~\newline ~\newline \satubaris{\PembimbingII}\newline NIDN. \PembimbingIInidn \\
\end{longtable}
\end{tabelpusat}\end{tanpagaris}
\begin{akhirsampul}\end{akhirsampul}
